\documentclass[11pt]{amsart}
\usepackage{amsmath,amssymb,amsthm}
\usepackage{booktabs}
\usepackage[colorlinks=true,allcolors=blue,linktoc=page]{hyperref}
\usepackage[capitalize]{cleveref}   % without the option nameinlink, only the number is a link

\theoremstyle{plain}
\newtheorem{theorem}{Theorem}[section]
\newtheorem{proposition}[theorem]{Proposition}
\newtheorem{lemma}[theorem]{Lemma}
\newtheorem{corollary}[theorem]{Corollary}
\theoremstyle{definition}
\newtheorem{definition}[theorem]{Definition}
\newtheorem{remark}[theorem]{Remark}
\newtheorem{question}[theorem]{Question}

\title[The second layer of the cyclotomic $\mathbb{Z}_7$-extension]{The class number of the second layer of the cyclotomic $\mathbb{Z}_7$-extension of $\mathbb{Q}$}
% \author{}   % to be filled in
\date{November 23, 2026 (draft)}

\emergencystretch=1.5em

\begin{document}

\begin{abstract}
We prove that the second layer of the cyclotomic $\mathbb{Z}_7$-extension of $\mathbb{Q}$, which is a cyclic field of degree $49$ and conductor $343$, has class number $1$.
The proof does not assume the generalized Riemann hypothesis.
It combines an upper bound for class numbers due to Miller, a divisibility result due to Masley, and explicit generators of the prime ideals above the $28381$ prime numbers up to $2.2 \cdot 10^{7}$ that split completely in this field.
\end{abstract}

\maketitle

\section{Introduction}

For each positive integer $m$, let $k_m$ denote the unique subfield of degree $m$ of the cyclotomic $\hat{\mathbb{Z}}$-extension of $\mathbb{Q}$.
If $p$ is a prime number and $n$ is a positive integer, then $k_{p^n}$ is the $n$-th layer of the cyclotomic $\mathbb{Z}_p$-extension of $\mathbb{Q}$.
This field is denoted by $\mathbb{B}_{p,n}$ in \cite{Mi15b}, and the field $k_{2^n}$ is often denoted by $\mathbb{B}_n$.
Let $h(k_m)$ denote the class number of $k_m$.

Weber \cite{We86} proved that $h(k_{2^n})$ is odd for every $n$ and that $h(k_{2^n}) = 1$ if $n \le 3$.
Whether $h(k_{2^n}) = 1$ holds for every $n$ is known as Weber's class number problem.
Coates \cite{Co12} asked whether $h(k_{p^n}) = 1$ holds for every prime number $p$ and every positive integer $n$.
At present, the equality $h(k_{p^n}) = 1$ is proved without any hypothesis for
\[
p^n \in \{2, 2^2, 2^3, 2^4, 2^5, 2^6,\ 3, 3^2, 3^3,\ 5, 5^2,\ 7,\ 11,\ 13,\ 17,\ 19\},
\]
and it is proved under the generalized Riemann hypothesis for
\[
p^n \in \{2^7,\ 3^4,\ 23,\ 29,\ 31\};
\]
these results are due to Weber, Cohn, Bauer, Masley, van der Linden, Coates, Liang, Mih\u{a}ilescu, and Miller; see \cite[Table 1, Theorem 1.8, and Table 2]{Mi15b} and \cite{Co60, Ba69, Ma78, vdL82, Co12, Mi14, Mi14b}.
For a fixed prime number $\ell$, the question whether $\ell$ divides $h(k_{p^n})$ for some $n$ has been studied by Horie and Horie \cite{HH09}, Fukuda and Komatsu \cite{FK11}, and Morisawa and Okazaki \cite{MO13}, among others.
An unconditional proof of $h(k_{23}) = 1$ by the same method is claimed in \cite{Ch26}; to our knowledge, it has not yet appeared in a refereed publication.

The purpose of this note is to add one field to the first list.

\begin{theorem} \label{thm:main}
The class number of $k_{7^2}$ is $1$.
\end{theorem}

The proof follows the method of Miller \cite{Mi15a, Mi15b}.
It has three steps.
First, an analytic inequality gives an upper bound for the class number, provided that sufficiently many prime ideals of degree one are known to be principal.
Second, a result of Masley shows that the class number is either $1$ or at least $197$.
Third, we exhibit generators of the prime ideals above every prime number $p \le 2.2 \cdot 10^7$ that splits completely in $k_{7^2}$.
With these prime ideals, the upper bound in the first step is smaller than $197$.

The field $k_{7^2}$ is within reach of this method for two reasons.
Its root discriminant is about $239.9$, which is small for a field of degree $49$, and the prime number $19$ splits completely in it.

The generators are too large to be printed.
They are available, together with the programs that verify every assertion of \cref{sec:cert,sec:bound}, at
\begin{center}
\url{https://github.com/uekijun46/Weber_class_number}
\end{center}

\section{The field} \label{sec:field}

Let $\zeta$ denote the complex number $e^{2\pi i/343}$.
The group $(\mathbb{Z}/343\mathbb{Z})^\times$ is cyclic of order $294 = 6 \cdot 49$.
Let $H$ denote its subgroup of order $6$, so that
\[
H = \{a \in (\mathbb{Z}/343\mathbb{Z})^\times \mid a^6 = 1\} = \{\pm 1, \pm 18, \pm 19\}.
\]
Let $K$ denote the subfield of $\mathbb{Q}(\zeta)$ fixed by $H$, where $a \in H$ acts by $\zeta \mapsto \zeta^a$.
Then $K = k_{7^2}$.
The field $K$ is totally real since $-1 \in H$, and it is cyclic of degree $49$ over $\mathbb{Q}$.
It is generated by the Gaussian period
\[
\theta = \sum_{a \in H} \zeta^a .
\]

\begin{lemma} \label{lem:field}
\begin{enumerate}
\item The discriminant of $K$ is $7^{138}$. The root discriminant of $K$ is $7^{138/49}$, and $\log 7^{138/49} = 5.4803183789\ldots$.
\item A prime number $p \ne 7$ splits completely in $K$ if and only if $p^6 \equiv 1 \pmod{343}$.
\item The ring of integers $O_K$ of $K$ is $K \cap \mathbb{Z}[\zeta]$.
\item For every $\alpha \in K$, we have $N_{\mathbb{Q}(\zeta)/\mathbb{Q}}(\alpha) = N_{K/\mathbb{Q}}(\alpha)^6$.
\end{enumerate}
\end{lemma}

\begin{proof}
(1) The group of Dirichlet characters associated with $K$ is cyclic of order $49$.
It consists of the trivial character, $6$ characters of conductor $7^2$, and $42$ characters of conductor $7^3$.
By the conductor-discriminant formula, the discriminant of $K$ is $7^{2 \cdot 6 + 3 \cdot 42} = 7^{138}$.

(2) The prime number $p$ splits completely in $K$ if and only if its Frobenius element in $\operatorname{Gal}(\mathbb{Q}(\zeta)/\mathbb{Q}) = (\mathbb{Z}/343\mathbb{Z})^\times$, which is the class of $p$, lies in $H$.

(3) The ring of integers of $\mathbb{Q}(\zeta)$ is $\mathbb{Z}[\zeta]$.

(4) This follows from $[\mathbb{Q}(\zeta):K] = 6$.
\end{proof}

\section{Three known results} \label{sec:known}

\subsection{An upper bound for class numbers}

Let $\gamma$ denote Euler's constant.
The following is \cite[Lemma 2.2]{Mi15b}; see also \cite[Lemma 3.1]{Mi15a}.

\begin{proposition}[Miller] \label{prop:miller}
Let $L$ be a totally real number field of degree $N$, and let $c$ be a positive real number.
Define a function $F$ on $\mathbb{R}$ by
\[
F(x) = \frac{e^{-(x/c)^2}}{\cosh(x/2)} .
\]
Let $S$ be a set of prime numbers each of which splits completely in $L$ into principal prime ideals.
Put
\[
I(c) = \int_0^\infty \frac{1 - F(x)}{2} \left( \frac{1}{\sinh(x/2)} + \frac{1}{\cosh(x/2)} \right) dx
\]
and
\[
B = \frac{\pi}{2} + \gamma + \log 8\pi - \log \operatorname{rd}(L) - I(c) + 2 \sum_{p \in S} \sum_{m=1}^{\infty} \frac{\log p}{p^{m/2}} F(m \log p),
\]
where $\operatorname{rd}(L)$ denotes the root discriminant of $L$.
If $B > 0$, then the class number $h$ of $L$ satisfies
\[
h < \frac{2c\sqrt{\pi}}{N B} .
\]
\end{proposition}

The set $S$ may be any set of prime numbers with the stated property.
Every term of the double sum is positive, so that a lower bound for $B$ is obtained if the sum is replaced by a partial sum.

\subsection{A divisibility result}

The following is the general form of \cite[Corollary 2.15]{Ma78}.
We include a proof for the convenience of the reader.

\begin{proposition}[Masley] \label{prop:masley}
Let $E/F$ be a cyclic extension of number fields of degree $n$, and let $q$ be a prime number that does not divide $n$.
Suppose that $q$ does not divide the class number of any field $L$ with $F \subseteq L \subsetneq E$.
Let $f$ denote the order of $q$ in $(\mathbb{Z}/n\mathbb{Z})^\times$.
If $q$ divides the class number of $E$, then $q^f$ divides the class number of $E$.
\end{proposition}

\begin{proof}
Let $G$ denote the Galois group of $E/F$, let $\sigma$ be a generator of $G$, and let $A$ denote the $q$-part of the ideal class group of $E$.
Since $q$ does not divide $n$, the ring $\mathbb{Z}_q[G]$ is the product of the rings $\mathbb{Z}_q[x]/(\Phi_e(x))$, where $e$ runs over the positive divisors of $n$ and $\Phi_e$ denotes the $e$-th cyclotomic polynomial.
Accordingly, $A$ is the direct sum of submodules $A_e$, where $\sigma$ acts on $A_e$ through $\mathbb{Z}_q[x]/(\Phi_e(x))$.

Let $e$ be a positive divisor of $n$ with $e < n$, and let $L$ denote the subfield of $E$ fixed by $\sigma^e$.
Then $F \subseteq L \subsetneq E$.
The element $\nu = \sum_{j=0}^{n/e-1} \sigma^{ej}$ of $\mathbb{Z}_q[G]$ acts on $A$ as the composite of the norm map to the $q$-part of the class group of $L$ and the map induced by extension of ideals.
The $q$-part of the class group of $L$ is trivial by assumption, so that $\nu$ acts on $A$ as zero.
On the other hand, $\sigma^e$ acts on $A_e$ as the identity, so that $\nu$ acts on $A_e$ as multiplication by $n/e$, which is a unit of $\mathbb{Z}_q$.
Hence $A_e = 0$.

Therefore $A = A_n$ is a module over $\mathbb{Z}_q[x]/(\Phi_n(x))$.
This ring is a product of copies of the ring of integers of the unramified extension of $\mathbb{Q}_q$ of degree $f$.
The order of every finite module over this ring is a power of $q^f$.
\end{proof}

\subsection{Two class numbers}

\begin{proposition} \label{prop:known}
\begin{enumerate}
\item The class number of $k_7$ is $1$.
\item The prime number $7$ does not divide the class number of $k_{7^2}$.
\end{enumerate}
\end{proposition}

\begin{proof}
(1) The field $k_7$ is the subfield of degree $7$ of $\mathbb{Q}(e^{2\pi i/49})$. Its class number was determined by Bauer \cite{Ba69}; see \cite[Table 1]{Mi15b}.
We also verified it with the function \texttt{bnfcertify} of PARI/GP \cite{PARI}, which does not assume any hypothesis.

(2) The extension $k_{7^2}/\mathbb{Q}$ is a Galois extension whose degree is a power of $7$, and $7$ is the only place of $\mathbb{Q}$ that ramifies in it.
Hence the assertion follows from a theorem of Iwasawa \cite{Iw56}; see \cite[Theorem 10.4]{Wa97}.
\end{proof}

\begin{corollary} \label{cor:197}
If $h(k_{7^2}) > 1$, then $h(k_{7^2}) \ge 197$.
\end{corollary}

\begin{proof}
Let $\ell$ be a prime number dividing $h(k_{7^2})$.
Then $\ell \ne 7$ by \cref{prop:known}.
The fields $L$ with $\mathbb{Q} \subseteq L \subsetneq k_{7^2}$ are $\mathbb{Q}$ and $k_7$, and both have class number $1$.
By \cref{prop:masley}, applied to $E = k_{7^2}$, $F = \mathbb{Q}$, and $n = 49$, the number $\ell^f$ divides $h(k_{7^2})$, where $f$ denotes the order of $\ell$ modulo $49$.

It remains to show that $\ell^f \ge 197$.
If $f = 1$, then $\ell \equiv 1 \pmod{49}$, and the least prime number with this property is $197$.
If $f \ge 2$ and $\ell^f < 197$, then $\ell \le 13$.
The orders of $2$, $3$, $5$, $11$, and $13$ modulo $49$ are $21$, $42$, $42$, $21$, and $14$, respectively, so that $\ell^f > 197$ in each of these cases.
\end{proof}

\section{Principal prime ideals} \label{sec:cert}

Let $S$ denote the set of prime numbers $p \le 2.2 \cdot 10^7$ with $p^6 \equiv 1 \pmod{343}$.
The set $S$ has $28381$ elements.
Its least elements are $19$, $1373$, $1697$, and $1733$, and its greatest element is $21997981$.
By \cref{lem:field}, every $p \in S$ splits completely in $K$.

\begin{proposition} \label{prop:cert}
For every $p \in S$, there exists an element $\alpha_p \in O_K$ with $|N_{K/\mathbb{Q}}(\alpha_p)| = p$.
Consequently, every prime ideal of $K$ above a prime number in $S$ is principal.
\end{proposition}

\begin{proof}
An element $\alpha_p$ is given for every $p \in S$ in the files accompanying this note, in the form
\[
\alpha_p = \frac{1}{D} \sum_{i=0}^{48} a_i \theta^i
\]
with integers $D \ge 1$ and $a_0, \ldots, a_{48}$.
The verification is carried out in the ring $\mathbb{Z}[x]/(\Phi_{343}(x)) = \mathbb{Z}[\zeta]$ with exact integer arithmetic.
The element $\sum_i a_i \theta^i \in \mathbb{Z}[\zeta]$ is computed, and it is checked that all of its coefficients with respect to the basis $1, \zeta, \ldots, \zeta^{293}$ are divisible by $D$.
Hence $\alpha_p \in K \cap \mathbb{Z}[\zeta] = O_K$ by \cref{lem:field}.
Then the resultant of $\Phi_{343}$ and the polynomial representing $\alpha_p$ is computed, and it is checked that its absolute value is $p^6$.
Hence $|N_{K/\mathbb{Q}}(\alpha_p)| = p$ by \cref{lem:field}.

Since the norm of the ideal $\alpha_p O_K$ is the prime number $p$, this ideal is a prime ideal above $p$.
The extension $K/\mathbb{Q}$ is Galois, so that every prime ideal of $K$ above $p$ is of the form $\sigma(\alpha_p) O_K$ with $\sigma \in \operatorname{Gal}(K/\mathbb{Q})$.
\end{proof}

\begin{remark} \label{rem:search}
The elements $\alpha_p$ were found with the functions \texttt{bnfinit} and \texttt{bnfisprincipal} of PARI/GP \cite{PARI}.
The function \texttt{bnfinit} assumes the generalized Riemann hypothesis.
This does not affect the proof of \cref{prop:cert}, since the proof only uses the elements $\alpha_p$ themselves, and these are verified independently.
The search took about $13$ seconds per prime number on average, which amounts to about $100$ hours on a single core, and the verification takes about $40$ minutes.
All $28381$ elements were also verified by a second program, which checks that $\alpha_p$ lies in the maximal order computed by PARI/GP and that the resultant of $\alpha_p$ and the minimal polynomial of $\theta$ has absolute value $p$.
\end{remark}

\section{The upper bound} \label{sec:bound}

We apply \cref{prop:miller} to $L = K$, $N = 49$, and the set $S$ of \cref{sec:cert}.
The hypothesis on $S$ holds by \cref{prop:cert}.

\begin{lemma} \label{lem:integral}
Let $c$, $\varepsilon$, and $M$ be positive real numbers with $\varepsilon < M$, and let $g$ denote the integrand in the definition of $I(c)$.
Then
\[
I(c) \le \int_\varepsilon^M g(x)\, dx + \left( \frac{1}{c^2} + \frac{1}{8} \right) \left( \frac{\varepsilon^2}{2} + \frac{\varepsilon^3}{6} \right) + 2 e^{-M/2} \left( 1 + \frac{1}{1 - e^{-M}} \right).
\]
\end{lemma}

\begin{proof}
Let $x$ be a positive real number.
From $\cosh(y) \le e^{y^2/2}$ and $e^{-u} \ge 1 - u$, we obtain
\[
F(x) \ge e^{-x^2/c^2 - x^2/8} \ge 1 - \frac{x^2}{c^2} - \frac{x^2}{8} .
\]
Together with $1/\sinh(x/2) \le 2/x$ and $1/\cosh(x/2) \le 1$, this gives
\[
g(x) \le \left( \frac{1}{c^2} + \frac{1}{8} \right) \left( x + \frac{x^2}{2} \right),
\]
and the integral of the right-hand side over the interval from $0$ to $\varepsilon$ is the second term of the assertion.

Suppose that $x \ge M$.
Then $0 \le 1 - F(x) \le 1$, $1/\sinh(x/2) = 2e^{-x/2}/(1 - e^{-x}) \le 2e^{-x/2}/(1 - e^{-M})$, and $1/\cosh(x/2) \le 2e^{-x/2}$.
Hence
\[
g(x) \le e^{-x/2} \left( 1 + \frac{1}{1 - e^{-M}} \right),
\]
and the integral of the right-hand side over the interval from $M$ to $\infty$ is the third term of the assertion.
\end{proof}

\begin{proposition} \label{prop:bound}
Let $c = 80$. Then
\begin{align*}
I(c) &\le 1.2651137825, \\
2 \sum_{p \in S} \sum_{m=1}^{8} \frac{\log p}{p^{m/2}} F(m \log p) &\ge 1.4092610134,
\end{align*}
and the number $B$ of \cref{prop:miller} satisfies $B > 0.03601227$.
Consequently, $h(k_{7^2}) < 160.7121$.
\end{proposition}

\begin{proof}
The two inequalities are obtained with ball arithmetic, using the library Arb \cite{Jo17}, which is now part of FLINT \cite{FLINT}.
For the first inequality, we use \cref{lem:integral} with $\varepsilon = 10^{-3}$ and $M = 200$, and the integral over the interval from $\varepsilon$ to $M$ is enclosed by the rigorous integration algorithm of \cite{Jo18}.
For the second inequality, each term is enclosed with a working precision of $200$ bits.

We have $\pi/2 + \gamma + \log 8\pi = 5.3721834192\ldots$.
By \cref{lem:field}, we have $\log \operatorname{rd}(K) = 5.4803183789\ldots$.
Hence
\begin{align*}
B &\ge 5.3721834192 - 5.4803183790 - 1.2651137825 + 1.4092610134 \\
&> 0.03601227,
\end{align*}
and \cref{prop:miller} gives
\[
h(k_{7^2}) < \frac{2 \cdot 80 \cdot \sqrt{\pi}}{49 \cdot 0.03601227} < 160.7121 . \qedhere
\]
\end{proof}

\begin{remark} \label{rem:values}
\Cref{tab:bounds} shows the bounds obtained in the same way for three values of $c$.
At $c = 80$, the inequality $B \ge 0.0293788$ suffices for $h < 197$.
The prime number $19$ alone contributes more than $0.6222$ to the sum over $S$ at $c = 80$, whereas the next element $1373$ of $S$ contributes less than $0.0209$.
Without the prime number $19$, the method would require a much larger set $S$.
\end{remark}

\begin{table}[h]
\centering
\begin{tabular}{rcccc}
\toprule
$c$ & $I(c) \le$ & sum $\ge$ & $B \ge$ & $h <$ \\
\midrule
$60$ & $1.2660224223$ & $1.3942300953$ & $0.0200727$ & $216.2490$ \\
$80$ & $1.2651137825$ & $1.4092610134$ & $0.0360122$ & $160.7121$ \\
$100$ & $1.2646927462$ & $1.4163260929$ & $0.0434983$ & $166.3167$ \\
\bottomrule
\end{tabular}
\caption{Bounds for $k_{7^2}$ obtained from \cref{prop:miller} with the set $S$. The third column is the double sum over $p \in S$ and $1 \le m \le 8$. Every entry is rounded in the safe direction, and the last column is computed from the value of $B$ before rounding.}
\label{tab:bounds}
\end{table}

\section{Proof of the theorem} \label{sec:proof}

\begin{proof}[Proof of \cref{thm:main}]
By \cref{prop:bound}, we have $h(k_{7^2}) < 160.7121$.
By \cref{cor:197}, the class number $h(k_{7^2})$ is either $1$ or at least $197$.
Hence $h(k_{7^2}) = 1$.
\end{proof}

\begin{remark} \label{rem:audit}
The numerical assertions of this note were checked in three ways.
The program accompanying this note verifies \cref{prop:cert,prop:bound} with exact integer arithmetic and ball arithmetic, and it does not use PARI/GP.
The elements $\alpha_p$ were verified again with PARI/GP, as described in \cref{rem:search}.
The bounds of \cref{tab:bounds} were recomputed with independent programs, in which the integral $I(c)$ is enclosed by Taylor models and the sum over $S$ by interval arithmetic.
\end{remark}

\begin{remark} \label{rem:next}
Under the generalized Riemann hypothesis, the analog of \cref{prop:miller} in \cite{Mi15b} requires only a small set $S$.
For the fields $k_{p^n}$ of larger degree, the difficulty then lies in finding generators of prime ideals.
Without the hypothesis, the fields $k_{29}$, $k_{3^4}$, and $k_{31}$ seem to be the next candidates for the present method, and they require sets $S$ with several million elements.
\end{remark}

\begin{thebibliography}{99}

\bibitem{Ba69} H.~Bauer, Numerische Bestimmung von Klassenzahlen reeller zyklischer Zahlk\"orper, J. Number Theory \textbf{1} (1969), 161--162.

\bibitem{Ch26} J.~Chakkour, \texttt{b23-class-number}, software repository, 2026, \url{https://github.com/johnchakkour/b23-class-number}.

\bibitem{Co12} J.~Coates, The enigmatic Tate--Shafarevich group, in: Fifth International Congress of Chinese Mathematicians, AMS/IP Stud. Adv. Math. \textbf{51}, Amer. Math. Soc., Providence, RI, 2012, pp.~43--50.

\bibitem{Co60} H.~Cohn, A numerical study of Weber's real class number of calculation. I, Numer. Math. \textbf{2} (1960), 347--362.

\bibitem{FK11} T.~Fukuda and K.~Komatsu, Weber's class number problem in the cyclotomic $\mathbb{Z}_2$-extension of $\mathbb{Q}$, III, Int. J. Number Theory \textbf{7} (2011), 1627--1635.

\bibitem{FLINT} The FLINT team, FLINT: Fast Library for Number Theory, \url{https://flintlib.org}.

\bibitem{HH09} K.~Horie and M.~Horie, The ideal class group of the $\mathbb{Z}_p$-extension over the rationals, Tohoku Math. J. \textbf{61} (2009), 551--570.

\bibitem{Iw56} K.~Iwasawa, A note on class numbers of algebraic number fields, Abh. Math. Sem. Univ. Hamburg \textbf{20} (1956), 257--258.

\bibitem{Jo17} F.~Johansson, Arb: efficient arbitrary-precision midpoint-radius interval arithmetic, IEEE Trans. Comput. \textbf{66} (2017), 1281--1292.

\bibitem{Jo18} F.~Johansson, Numerical integration in arbitrary-precision ball arithmetic, in: Mathematical Software -- ICMS 2018, Lecture Notes in Comput. Sci. \textbf{10931}, Springer, Cham, 2018, pp.~255--263; arXiv:1802.07942.
% CHECK: the page numbers

\bibitem{Ma78} J.~M.~Masley, Class numbers of real cyclic number fields with small conductor, Compositio Math. \textbf{37} (1978), 297--319.

\bibitem{Mi14} J.~C.~Miller, Class numbers of totally real fields and applications to the Weber class number problem, Acta Arith. \textbf{164} (2014), 381--397.

\bibitem{Mi14b} J.~C.~Miller, Class numbers of real cyclotomic fields of composite conductor, LMS J. Comput. Math. \textbf{17} (2014), suppl. A, 404--417.

\bibitem{Mi15a} J.~C.~Miller, Real cyclotomic fields of prime conductor and their class numbers, Math. Comp. \textbf{84} (2015), 2459--2469.

\bibitem{Mi15b} J.~C.~Miller, Class numbers in cyclotomic $\mathbb{Z}_p$-extensions, J. Number Theory \textbf{150} (2015), 47--73; arXiv:1410.2921.
% CHECK: the page numbers

\bibitem{MO13} T.~Morisawa and R.~Okazaki, Mahler measure and Weber's class number problem in the cyclotomic $\mathbb{Z}_p$-extension of $\mathbb{Q}$ for odd prime number $p$, Tohoku Math. J. \textbf{65} (2013), 253--272.

\bibitem{PARI} The PARI~Group, PARI/GP, Univ. Bordeaux, \url{https://pari.math.u-bordeaux.fr}.

\bibitem{vdL82} F.~J.~van der Linden, Class number computations of real abelian number fields, Math. Comp. \textbf{39} (1982), 693--707.

\bibitem{Wa97} L.~C.~Washington, Introduction to Cyclotomic Fields, 2nd ed., Grad. Texts in Math. \textbf{83}, Springer, New York, 1997.

\bibitem{We86} H.~Weber, Theorie der Abel'schen Zahlk\"orper, Acta Math. \textbf{8} (1886), 193--263.

\end{thebibliography}

\end{document}
